\SourceLeaf{200}
\[
AA+BB+CC,\quad AA'+BB'+CC',\quad A'A'+B'B'+C'C'
\]
\noindent Wenn den Endpunkten eines zweiten Elements auf der ersten Fläche die Werthe
\[
t,u\quad\text{und}\quad t+\delta t,u+\delta u
\]
\noindent entsprechen, so ist der Cosinus des Winkels, welchen dasselbe mit dem ersten Elemente macht,
\[
=\frac{\begin{gathered}(a\,\mathrm{d}t+a'\,\mathrm{d}u)(a\delta t+a'\delta u)+(b\,\mathrm{d}t+b'\,\mathrm{d}u)(b\delta t+b'\delta u)\\+(c\,\mathrm{d}t+c'\,\mathrm{d}u)(c\delta t+c'\delta u)\end{gathered}}{\begin{gathered}\BareRoot{((a\,\mathrm{d}t+a'\,\mathrm{d}u)^2+(b\,\mathrm{d}t+b'\,\mathrm{d}u)^2+(c\,\mathrm{d}t+c'\,\mathrm{d}u)^2)}\\\mathord{.}((a\delta t+a'\delta u)^2+(b\delta t+b'\delta u)^2+(c\delta t+c'\delta u)^2)\end{gathered}}
\]
\noindent und für den Cosinus des Winkels zwischen den correspondirenden Elementen auf der zweiten Fläche ergibt sich ein ganz ähnlicher Ausdruck, wenn nur \(a,b,c,a',b',c'\) in \(A,B,C,A',B',C'\) verwandelt werden. Offenbar werden beide Ausdrücke einander gleich, wenn die obige Proportionalität Statt findet, und die zweite Bedingung wird daher schon mit in der ersten begriffen, welches auch bei einigem Nachdenken von selbst klar ist.

Der analytische Ausdruck der Bedingung unserer Aufgabe ist demnach, dass
\[
\frac{AA+BB+CC}{aa+bb+cc}=\frac{AA'+BB'+CC'}{aa'+bb'+cc'}=\frac{A'A'+B'B'+C'C'}{a'a'+b'b'+c'c'}
\]
\noindent werden muss, welches eine endliche Function von \(t\) und \(u\) sein wird, die wir \(=mm\) setzen wollen. Es drückt dann \(m\) das Verhältniss aus, in welchem die Lineargrössen auf der ersten Fläche in ihrer Abbildung auf der zweiten vergrössert oder verkleinert werden (je nachdem \(m\) grösser oder kleiner ist als 1). Dieses Verhältniss wird, allgemein zu reden, nach den Stellen verschieden sein: in dem speciellen Falle, wo \(m\) constant ist, wird eine vollkommene Aehnlichkeit auch in den endlichen Theilen, und wenn überdiess \(m=1\) ist, wird eine vollkommene Gleichheit Statt finden, und die eine Fläche sich auf die andere abwickeln lassen.

\Article{5}
Indem wir Kürze halber
\[
\begin{aligned}
(aa+bb+cc)\mathrm{d}t^2+2(aa'+bb'+cc')\mathrm{d}t\mathord{.}\mathrm{d}u&\\
+(a'a'+b'b'+c'c')\mathrm{d}u^2&=\omega
\end{aligned}
\]
\noindent setzen, bemerken wir, dass die Differentialgleichung \(\omega=0\) \textit{zwei} Integrationen zulassen wird. Indem man nemlich das Trinomium \(\omega\) in zwei, in Beziehung auf
